\documentclass[12pt,a4paper]{article}

\usepackage{fontspec}
\usepackage{polyglossia}
\setmainlanguage{russian}
\setotherlanguage{english}
\setmainfont{Times New Roman}

\usepackage{amsmath,amssymb,amsthm,mathtools}
\usepackage{geometry}
\usepackage{enumitem}
\usepackage{array}
\usepackage{booktabs}
\usepackage{hyperref}
\usepackage{microtype}

\geometry{margin=2.5cm}
\setlength{\parindent}{1.25cm}
\setlength{\parskip}{0.35em}
\sloppy

\hypersetup{
    colorlinks=true,
    linkcolor=black,
    citecolor=black,
    urlcolor=blue,
    pdftitle={Решение проблем и уязвимостей Теории Единого Целого},
    pdfauthor={}
}

\theoremstyle{definition}
\newtheorem{definition}{Определение}[section]
\newtheorem{axiom}{Аксиома}[section]
\newtheorem{postulate}{Постулат}[section]

\theoremstyle{plain}
\newtheorem{theorem}{Теорема}[section]
\newtheorem{lemma}{Лемма}[section]
\newtheorem{proposition}{Предложение}[section]
\newtheorem{corollary}{Следствие}[section]

\theoremstyle{remark}
\newtheorem{remark}{Замечание}[section]

\newcommand{\Whole}{U}
\newcommand{\Parts}{\mathcal A}
\newcommand{\SubParts}{\mathcal B}
\newcommand{\Rel}{R}
\newcommand{\Values}{V}
\newcommand{\Vstruct}{\mathbf V}
\newcommand{\State}{S}
\newcommand{\States}{\mathcal S}
\newcommand{\PartOf}{\sqsubseteq}
\newcommand{\PropPart}{\sqsubset}
\newcommand{\comp}{\circ}
\newcommand{\Vord}{\preceq}
\newcommand{\norm}{\rho}
\newcommand{\ide}{e}
\newcommand{\Varr}{\operatorname{Var}}
\newcommand{\Stab}{\sigma}
\newcommand{\Disc}{\mathsf{Disc}}
\newcommand{\Cost}{\mathcal C}
\newcommand{\Time}{\mathsf{Time}}

\title{\textbf{Решение проблем и уязвимостей Теории Единого Целого}\\[0.5em]
\large Укрепление оснований: немножественная база, аксиоматизация значений,
вывод метрики, устранение скрытой арифметики, вывод кинетической энергии и
различающие предсказания}
\author{}
\date{}

\begin{document}

\maketitle
\tableofcontents
\newpage

\section*{Предварительное замечание}
\addcontentsline{toc}{section}{Предварительное замечание}

В критическом анализе предыдущих версий Теории Единого Целого (далее ТЕЦ) были
выделены шесть проблем, ставящих под сомнение внутреннюю строгость теории:

\begin{enumerate}
    \item[(3.1)] \textbf{Круговость в построении метрики.} Неравенство
    треугольника вводилось как требование к проекции, а не выводилось из аксиом.
    \item[(3.2)] \textbf{Скрытая арифметика.} Формула устойчивости и массы
    использовала сложение, модуль и деление над значениями отношений до введения
    числовой проекции.
    \item[(3.3)] \textbf{Заимствованная кинетическая энергия.} Выражение
    $\tfrac12 m v^2$ бралось из ньютоновской механики, а не выводилось.
    \item[(3.4)] \textbf{Теоретико-множественное основание.} Аксиомы записывались
    через $\in$ и $\subseteq$, несущие интуицию «пространства-контейнера».
    \item[(3.5)] \textbf{Отсутствие различающего предсказания.} Не было ни одного
    проверяемого отличия от существующих теорий.
    \item[(3.6)] \textbf{Незамкнутая структура значений.} Операции
    $\Vord,\comp,\norm$ вводились без аксиом.
\end{enumerate}

Настоящий документ решает все шесть проблем. Порядок изложения логический, а не
по нумерации: сначала строится немножественное основание (3.4), затем
аксиоматизируется структура значений (3.6), после чего метрика становится
теоремой (3.1), арифметика очищается (3.2), кинетическая энергия выводится (3.3),
и, наконец, из укреплённой теории извлекаются предсказания (3.5).

\section{Решение 3.4. Немножественное основание: мереология и решётка}

\subsection{Мотивировка}

Запись $A\subseteq\Whole$ и $X\in\Parts$ опирается на теорию множеств Цермело"--"Френкеля.
Множество в ней интуитивно мыслится как вместилище своих элементов, то есть несёт
скрытую пространственно-контейнерную семантику. Это противоречит цели ТЕЦ, в
которой пространство должно \emph{возникать}, а не предполагаться.

Решение состоит в замене теоретико-множественного основания на \emph{мереологию}
(исчисление «часть"--"целое») над \emph{ограниченным частичным порядком}. Отношение
«часть"--"целое» $\PartOf$ не является контейнерным: оно фиксирует лишь
структурный порядок, без точек, без вместимости, без расположения. Это
бесточечное (point-free) основание того же рода, что используется в теории локалей
и бесточечной топологии именно для того, чтобы строить геометрию без
предварительного допущения точек и пространства.

\subsection{Мереологические аксиомы}

Единственным первичным отношением объявляется бинарное отношение $\PartOf$
(«быть частью или совпадать») на области рассуждения $\mathfrak U$.

\begin{axiom}[Рефлексивность]
$x\PartOf x$ для всякого $x$.
\end{axiom}

\begin{axiom}[Антисимметрия]
Если $x\PartOf y$ и $y\PartOf x$, то $x=y$.
\end{axiom}

\begin{axiom}[Транзитивность]
Если $x\PartOf y$ и $y\PartOf z$, то $x\PartOf z$.
\end{axiom}

Тем самым $\PartOf$ есть частичный порядок.

\begin{axiom}[Единое Целое как наибольший элемент]
Существует единственный элемент $\Whole$, такой что $x\PartOf\Whole$ для всех $x$:
\[
\exists!\,\Whole\ \forall x:\ x\PartOf\Whole.
\]
\end{axiom}

\begin{axiom}[Реляционная полнота]
Для всякого непустого совокупления частей существует точная верхняя грань
относительно $\PartOf$. Таким образом $(\mathfrak U,\PartOf)$ есть ограниченная
полная верхняя полурешётка с наибольшим элементом $\Whole$.
\end{axiom}

\begin{definition}[Собственная часть и класс частей]
Часть называется собственной, если
\[
x\PropPart\Whole \;:\Longleftrightarrow\; x\PartOf\Whole\ \wedge\ x\neq\Whole.
\]
Классом частей называется
\[
\Parts:=\{x \mid x\PropPart\Whole\}.
\]
\end{definition}

\begin{definition}[Различимость без множеств]
Две части различимы, если они не совпадают как элементы порядка:
\[
\Disc(x,y)\;:\Longleftrightarrow\; x\neq y.
\]
\end{definition}

\begin{proposition}[Отсутствие контейнерной семантики]
В основании $(\mathfrak U,\PartOf,\Whole)$ нет ни отношения принадлежности $\in$,
ни понятия точки, ни понятия вместимости. Все отношения суть отношения порядка.
\end{proposition}

\begin{proof}
Сигнатура теории содержит единственный нелогический символ $\PartOf$, являющийся
частичным порядком. Понятия «элемента-точки», «расположения внутри» и «объёма» не
определимы средствами одного лишь порядка без дополнительной структуры. Эта
дополнительная структура (отношения $\Rel$, нормировка $\norm$) вводится лишь
ниже и порождает метрику как производное, а не как исходное.
\end{proof}

\begin{remark}[О метатеории]
Чтобы и понятие отображения $\Rel$ не вносило теоретико-множественной интуиции,
метатеорию допустимо принять не в виде ZFC, а в виде типовой теории: «отображение»
понимается как терм функционального типа $\Parts\times\Parts\to\Values$, а не как
множество пар. Это полностью устраняет возражение 3.4.
\end{remark}

\section{Решение 3.6. Аксиоматизация структуры значений отношений}

\subsection{Нормированный упорядоченный моноид значений}

Простого множества $\Values$ недостаточно: значения отношений нужно сравнивать,
компоновать и проецировать в число. Все эти операции должны быть аксиоматически
замкнуты, а единственным мостом к числам обязана быть нормировка $\norm$.

\begin{definition}[Структура значений отношений]
Структурой значений отношений называется набор
\[
\Vstruct=(\Values,\comp,\ide,\Vord,\norm),
\]
удовлетворяющий аксиомам V1--V7 ниже.
\end{definition}

\begin{axiom}[Моноид]
$(\Values,\comp,\ide)$ есть моноид: композиция $\comp$ ассоциативна, а $\ide$ "---" нейтральный элемент:
\[
(a\comp b)\comp c=a\comp(b\comp c),\qquad a\comp\ide=\ide\comp a=a.
\]
\end{axiom}

\begin{axiom}[Обратимость]
Каждый элемент обратим: существует операция $(\cdot)^{-1}$ с
\[
a\comp a^{-1}=a^{-1}\comp a=\ide,\qquad (a\comp b)^{-1}=b^{-1}\comp a^{-1}.
\]
Тем самым $(\Values,\comp,\ide)$ есть группа значений. (Если исходно задан лишь
моноид, рассматривается его обёртывающая группа Гротендика.)
\end{axiom}

\begin{axiom}[Порядок]
$\Vord$ есть частичный порядок на $\Values$.
\end{axiom}

\begin{axiom}[Монотонность композиции]
Порядок согласован с композицией слева и справа:
\[
a\Vord b\ \Longrightarrow\ a\comp c\Vord b\comp c\ \wedge\ c\comp a\Vord c\comp b.
\]
\end{axiom}

\begin{axiom}[Минимальность нейтрального элемента]
Нейтральный элемент является наименьшим:
\[
\ide\Vord a\quad\text{для всех }a\in\Values.
\]
\end{axiom}

\begin{axiom}[Нормировка-гомоморфизм]
Существует отображение $\norm:\Values\to\mathbb R_{\ge0}$ со свойствами:
\begin{align}
&\text{(V6a)}\quad \norm(\ide)=0;\\
&\text{(V6b)}\quad \norm(a\comp b)\le\norm(a)+\norm(b)\quad\text{(субаддитивность)};\\
&\text{(V6c)}\quad a\Vord b\ \Longrightarrow\ \norm(a)\le\norm(b)\quad\text{(монотонность)};\\
&\text{(V6d)}\quad \norm(a)=0\ \Longrightarrow\ a=\ide\quad\text{(положительная определённость)}.
\end{align}
\end{axiom}

\begin{axiom}[Симметрия нормы]
Нормировка инвариантна относительно обращения:
\[
\norm(a^{-1})=\norm(a).
\]
\end{axiom}

\begin{proposition}[$\norm$ есть абстрактная длина]
Из V1, V2, V6, V7 следует, что $\norm$ обладает всеми свойствами длины (нормы) на
группе: $\norm(a)\ge0$, $\norm(a)=0\iff a=\ide$, $\norm(a^{-1})=\norm(a)$ и
$\norm(a\comp b)\le\norm(a)+\norm(b)$.
\end{proposition}

\begin{proof}
Неотрицательность "--" по области значений $\norm$. Эквивалентность $\norm(a)=0\iff a=\ide$:
импликация $\Leftarrow$ есть V6a, импликация $\Rightarrow$ есть V6d. Симметрия "--" V7.
Субаддитивность "--" V6b.
\end{proof}

Таким образом возражение 3.6 снято: $\Vord,\comp,\norm$ полностью аксиоматизированы,
а структура $\Vstruct$ замкнута.

\subsection{Аксиомы отношений}

Отношение есть отображение $\Rel:\Parts\times\Parts\to\Values$, подчинённое
следующим аксиомам.

\begin{axiom}[Рефлексивность отношения]
$\Rel(A,A)=\ide$ для всякой части $A$.
\end{axiom}

\begin{axiom}[Различимость и нетривиальность]
\[
\Rel(A_i,A_j)=\ide\ \Longleftrightarrow\ A_i=A_j.
\]
Это реляционная форма тождества неразличимых: совпадение отношения с нейтральным
элементом равносильно совпадению частей.
\end{axiom}

\begin{axiom}[Обратимость отношения]
\[
\Rel(A_j,A_i)=\Rel(A_i,A_j)^{-1}.
\]
\end{axiom}

\begin{axiom}[Композиция отношений (субкомпозиция)]\label{ax:compose}
Для любых частей $A_i,A_j,A_k$:
\[
\Rel(A_i,A_k)\ \Vord\ \Rel(A_i,A_j)\comp\Rel(A_j,A_k).
\]
Прямое отношение не превосходит по порядку композиции отношения, проходящего через
любого посредника.
\end{axiom}

\begin{remark}
Аксиома~\ref{ax:compose} полностью внутренняя и реляционная: она не содержит ни
чисел, ни расстояний, ни арифметики. Именно она станет источником неравенства
треугольника. Содержательно она выражает принцип «прямое отношение есть наиболее
тесное»: посредник не может сделать связь теснее прямой.
\end{remark}

\section{Решение 3.1. Вывод метрики: неравенство треугольника как теорема}

\begin{definition}[Реляционное расстояние]
Реляционным расстоянием называется числовая проекция отношения:
\[
d(A_i,A_j):=\norm\bigl(\Rel(A_i,A_j)\bigr).
\]
\end{definition}

\begin{theorem}[Метрика выводится из аксиом]\label{thm:metric}
Отображение $d:\Parts\times\Parts\to\mathbb R_{\ge0}$ является метрикой. При этом
все четыре аксиомы метрики суть теоремы, а не требования.
\end{theorem}

\begin{proof}
\textbf{Неотрицательность.} $d(A_i,A_j)=\norm(\Rel(A_i,A_j))\ge0$ по области значений $\norm$.

\textbf{Тождество неразличимых.} По доказанному $\norm(a)=0\iff a=\ide$. Значит
\[
d(A_i,A_j)=0\iff \Rel(A_i,A_j)=\ide\iff A_i=A_j,
\]
где последняя эквивалентность есть аксиома различимости и нетривиальности.

\textbf{Симметрия.} По обратимости отношения и симметрии нормы:
\[
d(A_j,A_i)=\norm(\Rel(A_j,A_i))=\norm(\Rel(A_i,A_j)^{-1})=\norm(\Rel(A_i,A_j))=d(A_i,A_j).
\]

\textbf{Неравенство треугольника.} Применяя аксиому композиции, монотонность $\norm$
(V6c) и субаддитивность $\norm$ (V6b):
\[
d(A_i,A_k)=\norm(\Rel(A_i,A_k))
\ \le\ \norm\bigl(\Rel(A_i,A_j)\comp\Rel(A_j,A_k)\bigr)
\ \le\ \norm(\Rel(A_i,A_j))+\norm(\Rel(A_j,A_k))
=d(A_i,A_j)+d(A_j,A_k).
\]
Все четыре свойства установлены, следовательно $d$ "--" метрика.
\end{proof}

\begin{corollary}[Критерий метризуемости]
Реляционная система $(\Parts,\Rel)$ метризуема тогда и только тогда, когда значения
отношений образуют нормированную группу $\Vstruct$ с аксиомой композиции. Системы,
не допускающие такой нормировки $\norm$, неметризуемы; пространство в привычном
смысле есть частный случай реляционной структуры, а не её необходимое свойство.
\end{corollary}

\begin{remark}[Что именно изменилось]
В предыдущей версии неравенство треугольника постулировалось как условие на $F$.
Теперь оно \emph{выведено} из реляционной аксиомы композиции
(аксиома~\ref{ax:compose}) и гомоморфных свойств нормы. Круговость 3.1 устранена:
метрика рождается из структуры отношений, а не вкладывается в неё руками.
\end{remark}

\section{Решение 3.2. Устранение скрытой арифметики}

\subsection{Внутреннее определение изменения}

Прежняя запись $\sum_j|\Delta\Rel(A_i,A_j)|$ применяла модуль и сумму к значениям
отношений до всякой проекции. Это нелегитимно, пока $\Values\neq\mathbb R$. Исправление:
изменение определяется внутренне в группе значений, а арифметика выполняется только
\emph{после} применения $\norm$.

\begin{definition}[Реляционное приращение]
Пусть $\Rel_k(A_i,A_j)$ и $\Rel_{k+1}(A_i,A_j)$ "--" значения отношения в состояниях
$\State_k$ и $\State_{k+1}$. Реляционным приращением называется элемент группы значений
\[
\delta_{ij}:=\Rel_k(A_i,A_j)^{-1}\comp\Rel_{k+1}(A_i,A_j)\in\Values.
\]
Это единственный элемент, переводящий старое отношение в новое: $\Rel_{k+1}=\Rel_k\comp\delta_{ij}$.
\end{definition}

\begin{definition}[Вариация части]
Вариацией части $A_i$ между состояниями называется число, получаемое \emph{после}
нормировки каждого приращения:
\[
\Varr(A_i;\State_k,\State_{k+1}):=\sum_{j\neq i}\norm(\delta_{ij})\ \in\ \mathbb R_{\ge0}.
\]
Здесь суммирование уже законно: складываются вещественные числа $\norm(\delta_{ij})$,
а не значения отношений.
\end{definition}

\begin{definition}[Устойчивость и масса]
Устойчивость и масса определяются как функции вещественной вариации:
\[
\Stab(A_i):=\frac{1}{1+\Varr(A_i;\State_k,\State_{k+1})},
\qquad
m(A_i):=M\bigl(\Stab(A_i)\bigr),
\]
где $M:[0,1]\to\mathbb R_{\ge0}$ "--" фиксированная монотонная калибровочная функция.
\end{definition}

\begin{proposition}[Чистота арифметики]
В определениях вариации, устойчивости и массы нет ни одной арифметической операции
над значениями $\Values$. Единственный переход от $\Values$ к $\mathbb R$
осуществляется нормировкой $\norm$, после чего все операции выполняются в
$\mathbb R$.
\end{proposition}

\begin{proof}
Композиция $\comp$ и обращение $(\cdot)^{-1}$ "--" внутренние операции группы значений
(аксиомы V1, V2), а не арифметика. Приращение $\delta_{ij}$ получено только этими
операциями. Норма $\norm(\delta_{ij})\in\mathbb R_{\ge0}$. Сумма, прибавление единицы
и деление применяются исключительно к этим вещественным числам. Следовательно,
скрытая арифметика 3.2 устранена.
\end{proof}

\section{Решение 3.3. Вывод кинетической энергии без заимствования}

\subsection{Функционал стоимости изменения}

Множитель $\tfrac12$ и квадрат скорости ранее брались из ньютоновской механики.
Покажем, что они \emph{выводятся} как структура разложения Тейлора любого гладкого
функционала стоимости изменения в точке равновесия.

\begin{definition}[Реляционная скорость части]
Реляционной скоростью части $A_i$ называется темп нормированного изменения:
\[
v_i:=\frac{\Varr(A_i;\State_k,\State_{k+1})}{\Delta t},\qquad \Delta t=\tau(\State_k,\State_{k+1}).
\]
\end{definition}

\begin{definition}[Функционал стоимости]
Функционалом стоимости изменения называется отображение
$\Cost:\mathbb R\to\mathbb R_{\ge0}$, сопоставляющее реляционной скорости $v$
внутреннюю «стоимость» (действие) поддержания этого изменения.
\end{definition}

Вместо заимствования формулы постулируются лишь структурные свойства $\Cost$.

\begin{postulate}[Свойства стоимости изменения]\label{post:cost}
\begin{enumerate}
    \item[(E1)] $\Cost(0)=0$: покой (отсутствие изменения отношений) ничего не стоит.
    \item[(E2)] $\Cost$ достигает минимума при $v=0$: покой есть устойчивое
    равновесие, ибо устойчивость максимальна при отсутствии изменения. Отсюда
    $\Cost'(0)=0$.
    \item[(E3)] $\Cost$ дважды непрерывно дифференцируема в окрестности нуля.
    \item[(E4)] $\Cost$ чётна: $\Cost(v)=\Cost(-v)$. Стоимость изменения не зависит
    от знака (направления) изменения в пространстве значений (изотропия). Отсюда все
    нечётные производные в нуле равны нулю.
\end{enumerate}
\end{postulate}

\begin{theorem}[Квадратичность кинетической энергии]\label{thm:kinetic}
При выполнении E1--E4 для малых реляционных скоростей
\[
\Cost(v)=\tfrac12\,m\,v^2+O(v^4),\qquad m:=\Cost''(0)\ge0.
\]
Множитель $\tfrac12$ есть тейлоровский коэффициент $1/2!$, а масса $m$ есть кривизна
функционала стоимости в точке равновесия.
\end{theorem}

\begin{proof}
Разложение Тейлора $\Cost$ в нуле до четвёртого порядка:
\[
\Cost(v)=\Cost(0)+\Cost'(0)\,v+\tfrac12\Cost''(0)\,v^2+\tfrac16\Cost'''(0)\,v^3+O(v^4).
\]
По (E1) $\Cost(0)=0$. По (E2) $\Cost'(0)=0$. По (E4) чётности все нечётные
производные обращаются в нуль, в частности $\Cost'''(0)=0$. Следовательно
\[
\Cost(v)=\tfrac12\Cost''(0)\,v^2+O(v^4).
\]
Положим $m:=\Cost''(0)$. Поскольку $v=0$ "--" точка минимума, по условию второго
порядка $\Cost''(0)\ge0$, то есть $m\ge0$. Получаем $\Cost(v)=\tfrac12 m v^2+O(v^4)$.
\end{proof}

\begin{corollary}[Согласование массы как кривизны и массы как устойчивости]
Определение $m=\Cost''(0)$ согласовано с определением массы через устойчивость из
раздела~4. Большая кривизна $\Cost''(0)$ означает крутой рост стоимости изменения,
то есть сильное сопротивление изменению отношений, то есть высокую устойчивость
$\Stab(A_i)$ и большую массу $m(A_i)$. Обе линии "--" инерция как кривизна и инерция
как устойчивость "--" измеряют одно: сопротивление части изменению её отношений.
\end{corollary}

\begin{remark}[Что именно изменилось]
Ни $\tfrac12$, ни $v^2$ более не заимствуются. Они суть универсальная ведущая
асимптотика \emph{любого} гладкого чётного функционала стоимости с минимумом в
покое. Заимствование 3.3 устранено, а в качестве бонуса получен следующий
поправочный член $O(v^4)$, дающий проверяемое отличие (см. раздел~7).
\end{remark}

\section{Решение 3.5. Различающие предсказания}

Укреплённая теория порождает выводимые и проверяемые в принципе следствия. Приведём
три, прямо вытекающие из доказанного.

\subsection{Предсказание A. Поправки к кинетической энергии}

\begin{proposition}[Неньютоновские поправки]
Из теоремы~\ref{thm:kinetic} следует, что кинетическая энергия имеет вид
\[
K(v)=\tfrac12 m v^2+c_4 v^4+O(v^6),\qquad c_4=\tfrac{1}{24}\Cost^{(4)}(0).
\]
ТЕЦ предсказывает отклонение от ньютоновской зависимости $\tfrac12 m v^2$ при больших
реляционных скоростях, причём \emph{форма} ведущей поправки фиксирована (чётная,
степени $v^4$). Ньютоновская механика такой поправки не содержит; это отличие
проверяемо.
\end{proposition}

\subsection{Предсказание B. Инвариантная максимальная скорость}

\begin{theorem}[Из ограниченности нормы следует предельная скорость]
Пусть нормировка ограничена сверху, $\displaystyle\sup_{a\in\Values}\norm(a)=\norm_{\max}<\infty$
(отношения имеют максимальную интенсивность), и пусть длительность имеет минимальный
шаг $\tau_{\min}>0$. Тогда реляционная скорость ограничена инвариантной величиной
\[
v=\frac{\Delta d}{\Delta t}\le\frac{\norm_{\max}}{\tau_{\min}}=:c,
\]
не зависящей от выбора части и наблюдателя.
\end{theorem}

\begin{proof}
Изменение расстояния за один шаг ограничено диаметром значений:
$\Delta d=\norm(\delta)\le\norm_{\max}$. Длительность шага не меньше $\tau_{\min}$.
Отсюда $v=\Delta d/\Delta t\le\norm_{\max}/\tau_{\min}$. Величины $\norm_{\max}$ и
$\tau_{\min}$ суть структурные константы $\Vstruct$ и времени, одни и те же для всех
частей, поэтому $c$ инвариантна.
\end{proof}

\begin{remark}
Это реляционное происхождение инвариантной предельной скорости (аналога скорости
света). В ньютоновской схеме такого предела нет; ТЕЦ выводит его из ограниченности
нормы значений. Отличие принципиально и проверяемо.
\end{remark}

\subsection{Предсказание C. Квантование расстояния и длительности}

\begin{theorem}[Дискретный спектр расстояний]
Пусть порядок $(\Values,\Vord)$ фундирован и содержит атомы: существуют минимальные
нетривиальные значения $a_0$ (такие, что не существует $b$ с $\ide\prec b\prec a_0$),
а норма уважает атомы: $\norm(a_0)=\ell_0>0$ "--" минимальное положительное значение.
Тогда реляционное расстояние принимает значения из дискретного спектра, ограниченного
снизу величиной $\ell_0$, а длительность "--" из спектра с минимальным шагом $\tau_0>0$.
\end{theorem}

\begin{proof}
Любое нетривиальное значение $a\neq\ide$ мажорирует некоторый атом $a_0\Vord a$
(фундированность). По монотонности нормы $\norm(a)\ge\norm(a_0)=\ell_0$. Значит ненулевые
расстояния не меньше $\ell_0$, а спектр значений нормы на атомарно порождённых
элементах дискретен. Аналогичное рассуждение для функции длительности $\tau$ даёт
минимальный шаг $\tau_0$.
\end{proof}

\begin{remark}
ТЕЦ предсказывает минимальную длину $\ell_0$ и минимальную длительность $\tau_0$ "--"
дискретность пространства-времени на наименьшем масштабе. Это качественно отличает её
от континуальных полевых теорий и согласуется с ожиданиями феноменологии квантовой
гравитации.
\end{remark}

\subsection{Предсказание D. Размерная редукция}

\begin{remark}[Качественное следствие]
Размерность есть размерность вложения метрического пространства $(\Parts,d)$, то есть
число независимых инвариантов реляционной структуры. На наименьших масштабах, где
число различимых частей мало, число независимых инвариантов падает. ТЕЦ качественно
предсказывает уменьшение реляционной (спектральной) размерности на наименьших
масштабах "--" эффект, ожидаемый в ряде подходов к квантовой гравитации.
\end{remark}

\section{Сводка результатов}

\begin{center}
\renewcommand{\arraystretch}{1.4}
\begin{tabular}{>{\raggedright\arraybackslash}p{0.30\textwidth}>{\raggedright\arraybackslash}p{0.62\textwidth}}
\toprule
\textbf{Проблема} & \textbf{Решение} \\
\midrule
3.1. Круговость метрики &
Неравенство треугольника доказано (теорема~\ref{thm:metric}) из реляционной аксиомы
композиции~\ref{ax:compose} и гомоморфных свойств нормы; метрика выведена, а не
постулирована. \\
3.2. Скрытая арифметика &
Изменение определено внутренне через приращение $\delta_{ij}$ в группе значений;
норма $\norm$ применяется до арифметики; сумма и деление действуют лишь над числами. \\
3.3. Заимствованная энергия &
$\tfrac12 m v^2$ выведено (теорема~\ref{thm:kinetic}) как тейлоровская асимптотика
чётного функционала стоимости с минимумом в покое; масса $=\Cost''(0)$ "--" кривизна. \\
3.4. Множественное основание &
ZFC заменена мереологией над ограниченной решёткой $(\mathfrak U,\PartOf,\Whole)$;
ни $\in$, ни точек, ни контейнера; метатеория "--" типовая. \\
3.5. Нет предсказаний &
Получены поправки $O(v^4)$, инвариантная предельная скорость $c=\norm_{\max}/\tau_{\min}$,
квантование $\ell_0,\tau_0$ и размерная редукция. \\
3.6. Незамкнутая структура $\Values$ &
$\Values$ аксиоматизирована как нормированная упорядоченная группа
$\Vstruct=(\Values,\comp,\ide,\Vord,\norm)$ с полным набором аксиом V1--V7. \\
\bottomrule
\end{tabular}
\end{center}

\section*{Итог}
\addcontentsline{toc}{section}{Итог}

Все шесть проблем устранены не декларативно, а конструктивно. Логическая
последовательность укреплённой теории принимает вид
\[
(\mathfrak U,\PartOf,\Whole)
\ \to\ \Parts
\ \to\ \Disc
\ \to\ \Rel
\ \to\ \Vstruct=(\Values,\comp,\ide,\Vord,\norm)
\ \to\ \State
\ \to\ \Time
\ \to\ d
\ \to\ v,a
\ \to\ m,p,\mathfrak F,E,L,H,
\]
где каждая стрелка теперь снабжена аксиомами или доказанными теоремами, а единственным
мостом от структуры значений к числам служит нормировка $\norm$. В краткой форме:
\[
\boxed{\text{метрика, масса, энергия и стрела предельной скорости выводятся из}}
\]
\[
\boxed{\text{нормированной композиции отношений, а не вкладываются в теорию руками.}}
\]

\end{document}
